\section{n3. New CE theorem for every excellent normal local surface, arbitrary residue field and positive or mixed characteristic.}
\textit{Supports: G331, n1.}

NEW scoped theorem: every excellent normal local domain R of dimension TWO, in arbitrary equal or mixed characteristic and with arbitrary residue field, has a greatest continuous independently geometric C/D envelope for EVERY finite continuous signed obstacle on every primary section. The proof uses qualified G331 and own n1, and a dimension-specific published surface resolution theorem at its actual hypotheses. It does not assume higher-dimensional resolution, algebraizability, F-finiteness, CM or Q-Gorenstein hypotheses for arbitrary original R. It includes genuinely positive-characteristic and mixed-characteristic regular source CE in dimension two.

Step 1: actual regular primary models, proved rather than presumed cofinal. Given finitely many primary ideals, INCLUDE the gauge ideal a in this list and begin with their product blowup, which dominates their individual actual blowups by G331. Normalize it finitely by excellence. This is a normal projective integral surface, an isomorphism on U=Spec(R) minus the closed point; U is regular since R is normal of dimension two. Its singular locus is finite: excellence makes the regular locus open, normality gives regularity in codimension one, and a Noetherian closed zero-dimensional subset is finite. Every singular local completion is normal by excellence. Thus the exact condition(4) of Stacks54.14.5/Lipman holds. Its complete immediate proof, together with54.14.3--4 in the entire0BGJ unit, was read. It produces a regular resolution by finitely many normalized blowups. In that proof the local centres resolving this surface are its singular closed points, all over the marked point. Consequently the final regular surface Y is projective and unchanged on U. No analogous result in larger dimension is used.

Y carries an actual effective exceptional integral divisor E with O(-E) relatively ample and with positive coefficient on every exceptional component. On the initial normalized primary blowup the pullback tautological primary divisor is effective and its negative is ample, including under finite normalization. Each subsequent normalized blowup has its effective tautological divisor with relatively ample negative. Inductively combine this negative with a sufficiently large positive multiple of the pullback of the preceding negative ample divisor. Stacks29.38.7--8, whose whole immediate proofs were read, give ampleness of this combination. Pullbacks of effective Cartier divisors are effective here because every modification is birational integral and is generically outside their support. The initial primary divisor vanishes on the whole closed fibre; its pullback has positive order at every final exceptional component. Keeping it with positive multiplicity makes E have the asserted full support. One can add another positive multiple of it without spoiling ampleness, since its negative is generated by the original primary ideal. This also proves a nonempty exceptional fibre, even if R is regular, by including a proper primary blowup at the start.

Here is a full primary presentation argument, to ensure positivity takes place on an ACTUAL primary model. By proper birational normal-target descent, H\textasciicircum{}0(Y,O\_Y)=R (Stacks37.53.6, complete proof read); alternatively proper finiteness plus normal integral closure gives the same statement. L=O(-E) is ample on Y because the base is affine. Choose finitely many sections s\_i of powers of L whose nonvanishing loci Y\_(s\_i) are affine and cover Y. Make their degrees equal by powers, say n. Each affine Y\_(s\_i) is of finite type over R; choose finitely many R-algebra generators for its coordinate ring. The ample-section localization identity, proved in Stacks28.27.11, writes each generator as t/s\_i\textasciicircum{}k with t a section of L\textasciicircum{}(nk). Choose a common K for all these exponents; the sections s\_i\textasciicircum{}K and t s\_i\textasciicircum{}(K-k) of L\textasciicircum{}(nK) yield a projective-space map. In the s\_i\textasciicircum{}K chart their ratios contain all chosen affine generators, so Y\_(s\_i) maps as a closed subscheme of that affine chart. These charts cover the image; the map is an immersion. Y is proper over R and projective space is separated, so the immersion is proper and hence a closed immersion.

All the selected sections belong to R because L\textasciicircum{}(nK) subset O\_Y. Let J be their ideal in R. They generate L\textasciicircum{}(nK), so J O\_Y=L\textasciicircum{}(nK). The ideal is proper: a unit would be a section vanishing on the nonempty E. Off the marked point L is canonically O and generation shows J is the unit ideal. Thus J is proper primary. The projective-space map factors through Bl\_J(R), since J O\_Y is invertible. This blowup is the integral scheme Proj(R[Jt]) inside the same projective space. The factor is a closed immersion containing its generic point and indeed its dense U. A closed subscheme of an integral reduced scheme containing a dense open must have zero defining ideal. Hence Y IS Bl\_J(R). This proves precisely the regular primary refinements needed for finitely many ideal functions; no arbitrary-model cofinality is assumed. The finite section-generation/localization units used here were read with their whole immediate proofs.

Step 2: the finite degree problem. Write the exceptional curves of Y as E\_1,...,E\_r. Y is regular, so these integral Weil divisors are actual effective Cartier divisors (regular Noetherian local rings are factorial). No SNC or residue-field perfection is required. Write a O\_Y=O(-E\_a), E\_a=sum a\_i E\_i, a\_i positive integers, and let a signed model obstacle be f\_D, D=sum d\_i E\_i. Every finite real signed sum of primary gauges is of this form on the regular primary refinement above.

Define Q\_ij=deg\_(E\_i) O(E\_j), with degree over the residue field of R. For i!=j this is nonnegative: restriction gives an effective finite zero divisor, and its degree is the sum of local intersection lengths times residue degrees. The same length O\_(Y,x)/(f\_i,f\_j) computes Q\_ji, so Q is symmetric. Nonzero sections of any line bundle on a proper integral curve have nonnegative degree by the effective zero-divisor exact sequence. These arguments apply to nonsmooth integral curves and imperfect residue fields. All contracted proper integral curves on Y are among the E\_i. Thus an exceptional real divisor sum z\_i E\_i is curve-nef exactly when Qz>=0.

Put c\_i=min(d\_i,0). The feasible vectors are z<=c and Qz>=0. They are nonempty: if M is any nonnegative number with f\_D>=-M on X, the divisor -M E\_a is a negative effective nef barrier, since -E\_a is a pullback of the generated primary tautological bundle. Its coefficient vector b\_i=-M a\_i satisfies b<=c. Coefficientwise maximum preserves feasibility: at coordinate i choose whichever input has its ith coordinate equal to the maximum; the diagonal contribution is unchanged and every off-diagonal contribution increases because Q\_ij>=0. The subset b<=z<=c, Qz>=0 is compact. Maximize each of its finitely many coordinates and then take the finite coordinatewise maximum of these maximizers. The result z* is feasible and dominates the whole subset. Any original feasible z has max(z,b) in that subset, so z* is the GREATEST feasible vector of the entire problem.

Z*=sum z\_i* E\_i is an actual negative effective finite REAL Cartier presentation and has nonnegative degree on EVERY contracted integral curve. Y is an ACTUAL primary blowup by Step1. Therefore its evaluation H=f\_(Z*) belongs to G331's independent geometric core, hence continuous C. Coefficientwise Z*<=D implies H<=f\_D on EVERY support by actual Cartier effectivity. In particular this is actual geometric positivity, established before any envelope comparison.

Step 3: every competitor is controlled. Let q=-v(I)/s be a rational primary profile below f\_D, s a positive integer. On Y form F\_I=sum ord\_(E\_i)(I) E\_i, an actual effective Cartier divisor. Every generator of I is a global section of O(-F\_I): its height-one orders have the required lower bounds on the exceptional curves, and are nonnegative at horizontal primes; normal height-one intersection gives regularity after dividing by the Cartier generator. At the generic point of any E\_i, some ideal generator attains ord\_(E\_i)(I). Its resulting section of O(-F\_I)|\_(E\_i) is nonzero. Hence deg\_(E\_i) O(-F\_I)>=0 for EVERY i, by the curve argument above. Consequently A=-F\_I/s has Q-coefficients>=0 and is nonpositive.

The normalized divisorial points v\_i=ord\_(E\_i)/a\_i lie in X. Testing q<=f\_D there gives A\_i<=d\_i, so A is feasible and A<=Z*. Moreover I O\_Y subset O(-F\_I) gives q(v)<=f\_A(v) on every semivaluation, including nonzero support. To see this without an illicit fraction-field extension, pull both sheaves to the valuation ring of Frac(R/p), where p is that support. Properness gives its centre on Y; its generic point lies over U. The exceptional invertible lattice has a nonzero generic trivialization. Its local generator b and every ideal generator g have g/b regular at the centre, so v(g)>=v(b); a primary generator has finite value, and generators killed by p have infinite value and satisfy the same inequality. Thus q<=f\_A<=H EVERYWHERE. The unit case is identical with F\_I=0 when zero is a competitor.

N1 identifies the raw D envelope with the supremum of these rational competitors. It follows p(f\_D)<=H. The reverse follows because H itself is a C competitor. Hence p(f\_D)=H is the greatest C/D minorant, with actual core positivity. This is not norm averaging, a trace of a greatest minorant, or a claim that positivity follows from a raw supremum. No competing-model cofinality or arbitrary nef b-divisor pushforward is needed: the explicit ideal-section argument controls all competitors.

Step 4: all continuous obstacles and regularity. The signed real span of primary profiles is uniformly dense in C(X,R). For clarity the ideal identity proving the lattice condition is min(v(I)-v(J),v(K)-v(H))=v(IH+KJ)-v(JH), after common denominators; the span contains constants through v(a)=1. The truncated profiles v((g)+a\textasciicircum{}N)=min(v(g),N) separate every pair of semivaluations, including different supports. The compact lattice Stone-Weierstrass argument gives density, as also granted and proved in G331. Choose signed model f\_j uniformly converging to f. Their envelopes H\_j above are in C. N1's raw Lipschitz estimate makes H\_j uniformly Cauchy with limit p(f); closedness of C proves p(f) continuous and geometrically positive. Every canonical degree converges uniformly by n1. Exact primary-gauge changes and completion follow n1/G331. For decreasing continuous obstacles the actual D-envelope rule in n1 gives upper semicontinuous completed envelopes, including poles and bottom. No O(1/m) canonical rate or canonical FG for general surface model obstacles is claimed. A rational d yields a rational maximal coefficient vector by finite rational linear programming, but this does not identify its nef line bundle with a single ideal profile or prove semiampleness.

There is an explicit mixed-characteristic and positive-characteristic diagnostic instance. For a two-dimensional excellent regular local ring with parameters x,y and residue field k, blow up m=(x,y) and then the k-rational point y/x=0 of its exceptional curve. The affine point-blowup charts are regular: at a point of the exceptional curve their maximal ideals have the expected two generators and dimension two. The strict first curve E\_1 and second curve E\_2 have degree matrix Q=[[-2,1],[1,-1]]: the first self-degree decreases by one from -1, the new exceptional has degree -1, and their transverse intersection has length one. Pullback E\_m=E\_1+E\_2. The negative of 2E\_1+3E\_2 has strictly positive degrees on both curves. Stacks33.44.15 on the full closed fibre and30.21.4 over the local base, both with complete immediate proofs read, make it relatively ample. The primary presentation argument in Step1 then realizes this very two-blowup surface as an actual primary model after an ample power. For obstacle D=-E\_1-3E\_2, the greatest feasible vector is (-3/2,-3), with Qz*=(0,3/2). The ideal I=(x\textasciicircum{}2,y) becomes generated with divisor E\_1+2E\_2 on these two blowups: in the x chart it is x(x,y/x), whose second blowup principalizes the residual maximal ideal; at the other centres it is already invertible. Therefore the exact envelope is -3v(I)/2. This example is valid for k[[x,y]] in every characteristic and for Z\_p[[u]] with x=p,y=u; the latter is a complete excellent regular surface of mixed characteristic. It involves no Frobenius or characteristic-zero BFJ citation.

Own tiny exact program surface\_checks.py reports PASS on all225 coefficient obstacles d\_i=i/2, -10<=i<=4, for this two-curve matrix, and all925 pairs of their feasible vertices under coordinatewise maximum. It gives the example vector (-3/2,-3) and degrees (0,3/2). This is exhaustive finite rational linear-algebra evidence ONLY; it computes no ideals or semivaluations and does not prove the universal surface theorem. The universal proofs are Steps1--4 above.

The original general excellent task remains open in higher dimension and for unrestricted nonnormal conductor envelope descent. Surface source CE can be transported through the current413 uniquely lifting root maps or full orbit finite quotients at their verified hypotheses, but a finite cover or alteration without full-fibre greatest-minorant constancy supplies no general theorem. The curve clipping comparison and literal radial distinction remain n1. Products, energy and full finite-map operators remain separate. Artifacts: n3.txt, distinct expanded n3.tex, P12.3-surface-envelopes.tex, proof-companion.md, reports/surface-checks.json. Rerun: python surface\_checks.py; python build\_delivery.py; python check\_delivery.py. These are unverified mathematical draft proofs, not compilation or acceptance claims.

\subsection*{Expanded finite surface degree construction}
On the regular ACTUAL primary model, write
\[
 E_{\mathfrak a}=\sum_{i=1}^r a_i E_i,\qquad a_i>0,
 \qquad D=\sum_i d_i E_i,\qquad
 Q_{ij}=\deg_{E_i}\mathcal O_Y(E_j).
\]
For $i\ne j$, $Q_{ij}\ge0$ is the symmetric local intersection-length
sum. Every contracted integral curve is an $E_i$. If
$b_i=-Ma_i$ is the primary nef barrier and $c_i=\min(d_i,0)$, set
\[
 K=\{z\in\mathbb R^r:b\le z\le c,\ Qz\ge0\}.
\]
It is nonempty compact. For each $i$ choose a coordinate maximizer
$z^{(i)}\in K$. If $u=\max(z,w)$, choose $z_i=u_i$; then
\[
 (Qu)_i=Q_{ii}z_i+\sum_{j\ne i}Q_{ij}u_j
 \ge (Qz)_i\ge0.
\]
Thus $z^*=\max_i z^{(i)}\in K$ is its greatest vector. Replacing any
unbounded feasible vector $z$ by $\max(z,b)$ proves it is greatest
in the full problem as well. The actual negative effective divisor
$Z^*=\sum_i z_i^*E_i$ is curve-nef before the envelope is identified.

For every rational competitor $q_I=-v(I)/s$, put
\[
 F_I=\sum_i\operatorname{ord}_{E_i}(I)E_i,\qquad A=-F_I/s.
\]
The inclusion $I\mathcal O_Y\subset\mathcal O_Y(-F_I)$ and a generator
attaining the generic minimum on each $E_i$ give a nonzero section of
$\mathcal O_Y(-F_I)|_{E_i}$. Its degree is nonnegative. Hence $QA\ge0$.
Testing $q_I\le f_D$ at $\operatorname{ord}_{E_i}/a_i$ yields
$A_i\le d_i$, and the valued quotient-field lattice proves
\[
 q_I(v)\le f_A(v)\le f_{Z^*}(v)\qquad\text{for EVERY support}.
\]
The rational-competitor identity in n1 therefore gives
$p(f_D)=f_{Z^*}\in C_{\rm geom}$. Uniform signed-model density and
$\|p(f)-p(g)\|\le\|f-g\|$ give the theorem for every continuous signed
obstacle. There is no claimed canonical approximation rate.

The two-blowup computation is
\[
 Q=\begin{pmatrix}-2&1\\1&-1\end{pmatrix},\quad
 d=\binom{-1}{-3},\quad z^*=\binom{-3/2}{-3},\quad
 Qz^*=\binom0{3/2}.
\]
With $I=(x^2,y)$ and $F_I=E_1+2E_2$, the actual envelope is
$p(f_D)=-\tfrac32 v(I)$, also on $\mathbb Z_p[[u]]$ for $x=p,y=u$.
